% !TEX program = xelatex
\documentclass[UTF8,a4paper,12pt]{ctexart}

\usepackage{amsmath,amssymb}
\usepackage{booktabs}
\usepackage{caption}
\usepackage{float}
\usepackage{graphicx}
\usepackage{geometry}
\usepackage{indentfirst}
\usepackage{siunitx}

\geometry{left=2.6cm,right=2.6cm,top=2.5cm,bottom=2.5cm}
\captionsetup{font=small,labelsep=quad}
\setlength{\parindent}{2em}
\setlength{\parskip}{0pt}
\linespread{1.35}
\sisetup{detect-all}

% 若总论文已定义该命令，本行不会覆盖原定义。
\providecommand{\QThreeDryHours}{57.5833}

% 图件尚未放入目录时仍可正常编译；补齐同名图片后会自动替换占位框。
\newcommand{\safeincludegraphics}[2][]{%
  \IfFileExists{#2}{%
    \includegraphics[#1]{#2}%
  }{%
    \fbox{\parbox[c][5.4cm][c]{0.82\textwidth}{\centering
      \textbf{图件待补充}\\[0.6em]
      请将 \texttt{\detokenize{#2}} 放入本文件所在目录，\\
      或在总论文中修改为实际图片路径。}}%
  }%
}

\begin{document}

\section{模型求解与结果分析}

\subsection{问题3：固定半径下的干燥终点}

\subsubsection{问题分析与终止条件}

问题3是在物料半径保持不变的条件下，求出其全域水分浓度均降至
$0.15$ 以下的最早时刻。本题不用再建立新的传热传质模型，而是直接沿用
问题2所得到的温度场--水分场联立系统，并将计算时域延长至干燥过程结束。
需要注意的是，表面节点较早达到给定阈值并不能说明物料已经完全干燥；
只有径向区域内水分浓度的最大值低于阈值时，才可判定整个物料满足要求。

定义全场最大水分浓度
\begin{equation}
M(t)=\max_{0\leq r\leq R}C(r,t),
\label{eq:q3_max_concentration}
\end{equation}
并构造事件函数
\begin{equation}
g(t)=M(t)-0.15.
\label{eq:q3_event}
\end{equation}
由于干燥阶段的控制量为全域最大值，理论临界时刻可表示为
\begin{equation}
t_{\mathrm{c}}
=\inf\left\{t>0\,\middle|\,M(t)\leq 0.15\right\},
\qquad g\!\left(t_{\mathrm{c}}\right)=0.
\label{eq:q3_continuous_time}
\end{equation}
(1)式中，$R$为物料半径，$C(r,t)$为位置$r$、时刻$t$处的水分浓度。

\subsubsection{数值求解方法}

在问题2的求解过程中，我们已经得到了物料内部温度和水分浓度随时间的
变化情况。因此，对于问题3，只需保持物料半径不变，将原有模型的求解时间
适当延长，直到物料各处的水分浓度均低于$0.15$即可。考虑到程序每隔
$60\ \mathrm{s}$输出一次结果，理论上的临界时刻不一定恰好落在输出节点上，
所以可以先利用事件函数确定其准确位置，再从输出结果中选取一个满足要求的安全
时刻。具体求解步骤如下：

\noindent\textbf{Step1：}输入问题2中已经确定的模型参数、初始条件和边界条件，
保持物料半径$R$不变，并将径向区域划分为801个节点。

\noindent\textbf{Step2：}用问题2中的温度场--水分场联立方程进行计算。分别求出各径向节点处每时刻的水分浓度，并从中选出最大值$M(t)$。

\noindent\textbf{Step3：}判断$M(t)$是否降至$0.15$。若事件函数$g(t)$由正值
变为非正值，则说明全域最大水分浓度在这一时间段内达到规定要求，即可进一步在
相邻两个计算时刻之间求解$g(t)=0$，从而得到理论临界时刻
$t_{\mathrm{c}}$；若不满足，则需继续延长时间进行计算。

\noindent\textbf{Step4：}由于题目结果在$60\ \mathrm{s}$时间网格上给出，还需在
$t_n=60n$处进行再次判断。若某一输出时刻满足$M(t_n)<0.15$，则将其作为最终的
安全达标时刻；否则，需继续考察下一个输出时刻。即
\begin{equation}
t_3=\min\left\{60n\,\middle|\,
\max_{0\leq j\leq800}C(r_j,60n)<0.15\right\}.
\label{eq:q3_safe_time}
\end{equation}

\noindent\textbf{Step5：}输出最终达标的时刻，并记录物料表面和中心位置达到
$0.15$的时间。最后结合温度场和水分场的变化图，对干燥终点进行分析。

通过上述步骤，既可以得到水分浓度恰好达到$0.15$的理论时刻，又能保证所取
的离散输出时刻已经满足全域水分浓度低于$0.15$的要求。

\subsubsection{计算结果与分析}

通过数值计算可得，事件函数首次到达零点的连续临界时刻为
\begin{equation}
t_{\mathrm{c}}=57.5685\ \mathrm{h}
=207246.6\ \mathrm{s}.
\label{eq:q3_critical_result}
\end{equation}
在$60\ \mathrm{s}$输出网格上继续进行检验，首个满足全域水分浓度严格低于
$0.15$的安全时刻为
\begin{equation}
\boxed{t_3=207300\ \mathrm{s}
=\QThreeDryHours\ \mathrm{h}
\approx2.40\ \mathrm{d}}.
\label{eq:q3_final_result}
\end{equation}
此时全域最大水分浓度为$0.149984$，比规定阈值低$1.6\times10^{-5}$。
两级判停结果如表\ref{tab:q3_results}所示。

\begin{table}[H]
\centering
\caption{问题3干燥终点计算结果}
\label{tab:q3_results}
\begin{tabular}{cccc}
\toprule
判定方式 & 达标时刻/s & 达标时刻/h & 全域最大水分浓度 \\
\midrule
连续事件定位 & 207246.6 & 57.5685 & 0.150000 \\
$60\ \mathrm{s}$网格复核 & 207300 & 57.5833 & 0.149984 \\
\bottomrule
\end{tabular}
\end{table}

由表1可以看出，两种方法所得时刻仅相差$53.4\ \mathrm{s}$，小于程序的一个
输出步长，因而选取$207300\ \mathrm{s}$作为最终结果是合理的。此外，通过对物料
不同位置的水分浓度进行比较可以发现，物料表面处的水分浓度在$12.55\ \mathrm{h}$左右已经降至
$0.15$，而中心处的水分浓度约在$57.58\ \mathrm{h}$才降至这一数值，前后相差
$45.03\ \mathrm{h}$。

这是由于物料表面的水分可以直接向外界扩散，而中心处的水分
需要先从物料内部逐渐向表面迁移，这一过程所经过的距离较长，因而下降速度较慢。因此，
整个物料何时能够达到干燥要求，这主要取决于中心位置的水分浓度。如果只根据
表面位置进行判断，所得干燥时间将会偏小。

全域最大水分浓度随时间的变化曲线如图\ref{fig:drying}所示。由图1可以看出，
随着干燥时间的增加，$M(t)$不断减小，并在$57.5685\ \mathrm{h}$附近降至
$0.15$。此时继续计算至下一个$60\ \mathrm{s}$输出时刻，得到的全域最大水分浓度为
$0.149984$，已经满足题目要求。

\begin{figure}[H]
\centering
\safeincludegraphics[width=0.82\textwidth]{drying_threshold.png}
\caption{全域最大水分浓度随时间的演化及$60\ \mathrm{s}$安全达标点定位}
\label{fig:drying}
\end{figure}

图\ref{fig:q3_heatmaps} 给出了物料达到干燥要求前后温度场和水分场的变化情况。从图中可以看出，干燥开始后，
物料表面处的水分浓度下降较快，而中心处下降相对较慢。在干燥后期，物料内部
仍存在一定的水分浓度差异，中心处的水分浓度最后达到规定值，这与前面所得表面和中心
位置的达标时间相符。

\begin{figure}[H]
\centering
\safeincludegraphics[width=0.92\textwidth]{question3_heatmaps.png}
\caption{问题3直至达标后$0.5\ \mathrm{h}$的温度场与水分场二维时空演化热力图}
\label{fig:q3_heatmaps}
\end{figure}

根据上述计算，在固定半径条件下，物料达到干燥要求的时间为
$207300\ \mathrm{s}$。由于中心处的水分浓度下降最慢，因而只有当物料各位置
水分浓度的最大值低于$0.15$时，才说明整个物料已经达到要求。如果只考察物料表面
位置的水分浓度，则会过早结束计算，所得到的干燥时间也会小于实际所需时间。

\end{document}
